% Generated by generate_ieee14_gfm_lock_comparison. Do not edit.
% Classification: ASSUMED_DIAGNOSTIC reporting over production runs.

\begin{table}[H]\centering\footnotesize
\caption{Arm outcomes on the 250-s IEEE 14-bus chronology. Identical case, dispatch, event schedule, solver and step control; the arms differ only in the converter control policy. \textbf{--} means the quantity is not defined for that arm, never zero.}
\begin{tabular}{lrrrrr}\toprule
quantity & adaptive & pinned 1 GFM & pinned 2 GFM & pinned 4 GFM & locked GFL \\ \midrule
horizon reached [s] & 250.000 & 250.000 & 250.000 & 25.491 & 20.000 \\
converged & 1 & 1 & 1 & 0 & 0 \\
refusal identifier & -- & -- & -- & \texttt{adaptiveDtMin} & \texttt{noVoltageFormingSource} \\
reclose status & SUCCESS & SUCCESS & SUCCESS & NOT\_REQUESTED & NOT\_REQUESTED \\
reclose instant [s] & 159.2397 & 151.4244 & 149.9191 & -- & -- \\
grid-forming units, max & 4 & 1 & 2 & 4 & 0 \\
grid-forming units, final & 1 & 1 & 2 & 4 & 0 \\
supervisor commits applied & 4 & 0 & 0 & 0 & 0 \\
terminal $f_{COI}$ [Hz] & 60.000000 & 60.000000 & 60.000000 & -- & -- \\
accepted samples & 4990 & 3232 & 3183 & 819 & 43 \\
rejected steps & 135 & 542 & 321 & 449 & 0 \\
\bottomrule\end{tabular}\end{table}

\begin{table}[H]\centering\footnotesize
\caption{Frequency excursion and deepest network voltage in the 15\,s after each disturbance. \textbf{--} means the arm never reached that event.}
\begin{tabular}{lrrrrr}\toprule
window & adaptive & pinned 1 GFM & pinned 2 GFM & pinned 4 GFM & locked GFL \\ \midrule
\multicolumn{6}{l}{\emph{$\max|f_{COI}-60|$ [Hz]}} \\
\quad SG trip ($t=20$\,s) & 0.7587 & 0.5312 & 0.3489 & 0.9898 & -- \\
\quad load step ($t=50$\,s) & 1.0035 & 1.4478 & 0.7839 & -- & -- \\
\quad fault ($t=85$\,s) & 1.1386 & 1.4979 & 0.8255 & -- & -- \\
\quad line trip ($t=110$\,s) & 0.3764 & 1.2325 & 0.7135 & -- & -- \\
\quad restore/reclose ($t=145$\,s) & 1.1612 & 1.1736 & 0.6634 & -- & -- \\
\multicolumn{6}{l}{\emph{$\min|V|$ [pu]}} \\
\quad SG trip ($t=20$\,s) & 0.5508 & 0.5508 & 0.5508 & 0.5508 & -- \\
\quad load step ($t=50$\,s) & 0.8128 & 0.8128 & 0.8011 & -- & -- \\
\quad fault ($t=85$\,s) & 0.0369 & 0.0359 & 0.0364 & -- & -- \\
\quad line trip ($t=110$\,s) & 0.9610 & 0.9390 & 0.9502 & -- & -- \\
\quad restore/reclose ($t=145$\,s) & 0.9562 & 0.9396 & 0.9246 & -- & -- \\
\multicolumn{6}{l}{\emph{settled $\overline{|f_{COI}-60|}$ [Hz], last 30\% of the window}} \\
\quad SG trip ($t=20$\,s) & 0.0041 & 0.0113 & 0.0001 & -- & -- \\
\quad load step ($t=50$\,s) & 0.3753 & 1.2214 & 0.7123 & -- & -- \\
\quad fault ($t=85$\,s) & 0.3753 & 1.2296 & 0.7126 & -- & -- \\
\quad line trip ($t=110$\,s) & 0.3503 & 1.1426 & 0.6422 & -- & -- \\
\quad restore/reclose ($t=145$\,s) & 0.0130 & 0.0210 & 0.0209 & -- & -- \\
\bottomrule\end{tabular}\end{table}

\begin{table}[H]\centering\small
\caption{Agreement on the common window $[0,\,20)$\,s, before the event at which the policies first differ. This is what makes the later difference attributable to the policy alone.}
\begin{tabular}{lrrrr}\toprule
pair & samples & $\max|\Delta x|$ & $\max|\Delta y|$ & bit-identical \\ \midrule
adaptive vs pinned\_gfm1 & 42 & $0$ & $0$ & yes \\
adaptive vs pinned\_gfm2 & 42 & $0$ & $0$ & yes \\
adaptive vs pinned\_gfm4 & 42 & $0$ & $0$ & yes \\
adaptive vs locked\_gfl & 42 & $0$ & $0$ & yes \\
\bottomrule\end{tabular}\end{table}

\paragraph{Scope of this comparison.}
Every arm uses the same network, the same dispatch, the same event schedule, the same solver and the same step control, and the runner asserts that the realised option structs differ only in the fields the arm declares. The arms are bit-identical on $[0,\,20)$\,s, so the difference after that instant is attributable to the converter control policy and to nothing else.
\emph{What the locked grid-following arm establishes.} With every converter locked in grid-following the synchronous machine is the only voltage-forming source, so opening its breaker leaves an energised island with no angle reference. The per-island admissibility check refuses the transaction before any Newton solve and the trajectory ends at the event-left sample. On this resource mix, promoting at least one converter to grid-forming is therefore a \emph{necessary condition} for a post-trip trajectory to exist at all --- not an improvement to one.
\emph{What the pinned single-unit arm establishes.} One pinned grid-forming unit is enough to keep the island well posed, and that arm also reaches the horizon. The comparison is therefore quantitative rather than binary. Over the 5 disturbance windows both arms reached, the adaptive policy gives the smaller peak frequency excursion in 4 and the larger one in 1.
The larger effect is not the peak but the settled deviation. In the fault window the adaptive island holds $\overline{|f_{COI}-60|}=0.3753$\,Hz against 1.2296\,Hz for the single pinned unit, a factor of 3.28. The mechanism is droop sharing: the supervisor had promoted 4 units by then, so each carries a fraction of the imbalance and the common steady-state deviation is correspondingly smaller. That is the operative difference between the two policies on this study.
The adaptive advantage is largest at the disturbances that follow an earlier promotion, because the supervisor had already left the island with more grid-forming capacity committed. At the SG trip itself the pinned arm is the better of the two, and that is reported as found: the adaptive arm pays a small transient for its own promotion sequence.
\emph{What this comparison does \textbf{not} establish.}
\begin{itemize}\itemsep0pt
\item No percentage improvement against the locked arm over $[20,\,250]$\,s. That arm produces no trajectory there, and none is fabricated.
\item Nothing about a system that has another voltage-forming source. An all-synchronous arm is deferred and is a separate study.
\item Nothing about the choice of \emph{which} converters to promote against a different selector. The comparison is switching versus not switching.
\item Nothing numerical. Residuals, subdivision depth and rejected-step counts are not comparable across arms of different length.
\item Every reference-AGSI sub-index beyond $J_V$ and $J_f$ is \textsf{ASSUMED\_DIAGNOSTIC}; it entered no gate and supports no readiness claim.
\end{itemize}
